\section{Time regularity and a differentiable speed-shell residual}
\label{st:section}

\paragraph{Scope and status.}
These are author-checked component derivations, awaiting independent
mathematical audit. They resolve a temporal regularity obstacle for a
\emph{finite, regularized} shell residual. They do not bound the critical
clock at a possible singular endpoint. The natural variable is
$V=|w|^{1/2}w$, not $w$ itself or its unit direction. The underlying
monotone-map geometry is standard nonlinear variational geometry; no
novelty claim is made. All constants and limiting arguments used below
are supplied here. The Riesz bounds are the same classical
Calder\'on--Zygmund inputs used earlier, not additional PDE estimates.

\subsection{A natural-distance estimate with gradient freedom}

For $z\in\mathbb R^3$ set $A(z)=|z|z$ and $J(z)=|z|^{1/2}z$.
Let $\mathcal G_3$ be the closed gradient subspace and $w(f)$ the cubic
minimizer of $f+\mathcal G_3$, for any $f\in L^3$.

\begin{theorem}[Natural-variable stability]\label{st:stability}
Let $w_i=w(f_i)$, $\rho_i=|w_i|$, and $V_i=J(w_i)$, for
$f_0,f_1\in L^3(\mathbb R^3;\mathbb R^3)$. For every $g\in\mathcal G_3$,
\begin{equation}\label{st:difference}
 \|V_1-V_0\|_2^2
 \le\frac98\int(\rho_0+\rho_1)|f_1-f_0+g|^2.
\end{equation}
In particular $f\mapsto J(w(f))$ is Lipschitz from every bounded subset
of $L^3$ into $L^2$, with squared Lipschitz constant at most $9R/4$
on $\{\|f\|_3\le R\}$.
\end{theorem}
\begin{proof}
Both maps are $C^1$ at zero, with zero derivative there. For $z\ne0$,
writing $r=|z|$ and $e=z/r$, their derivatives are
\[
 DA(z)=r(I+e\otimes e),\qquad
 DJ(z)=r^{1/2}(I+\tfrac12e\otimes e).
\]
Thus, for every vector $d$,
\begin{equation}\label{st:matrix}
 |DJ(z)d|^2\le\tfrac98\,d^{\mathsf T}DA(z)d.
\end{equation}
The ratios in the radial and transverse directions are respectively
$9/8$ and $1$; at zero both sides vanish. Along the segment from $a$ to
$b$, Cauchy--Schwarz in its parameter and \eqref{st:matrix} give
\[
 |J(b)-J(a)|^2\le\tfrac98
 (A(b)-A(a))\cdot(b-a).
\]
Set $d=w_1-w_0$ and
$H(x)=\int_0^1DA(w_0+s d)\,ds$. This is a measurable positive
semidefinite matrix, and
\[
 A(w_1)-A(w_0)=Hd,\qquad 0\le H\le(\rho_0+\rho_1)I.
\]
The last inequality follows from $DA(z)\le2|z|I$ and the convexity
bound $|w_0+s d|\le(1-s)\rho_0+s\rho_1$.
The Euler--Lagrange condition says that both $A(w_i)$ annihilate
$\mathcal G_3$. Since $d-(f_1-f_0+g)\in\mathcal G_3$, putting
$e_0=f_1-f_0+g$ yields
\[
 E:=\int d^{\mathsf T}Hd
   =\int d^{\mathsf T}He_0
   \le E^{1/2}\left(\int e_0^{\mathsf T}He_0\right)^{1/2}.
\]
Every integral is finite by H\"older in $L^3,L^3,L^3$.
Cauchy--Schwarz is applied to $H^{1/2}d,H^{1/2}e_0$ and does not
require invertibility of $H$. If $E=0$ the conclusion is immediate;
otherwise division gives $E\le\int(\rho_0+\rho_1)|e_0|^2$.
Integrate the pointwise natural-distance inequality to prove
\eqref{st:difference}. Finally $\|w(f_i)\|_3\le\|f_i\|_3$ and
H\"older give the stated Lipschitz constant by taking $g=0$.
\end{proof}

\begin{theorem}[Pressure-free local time derivative]\label{st:time}
On every compact interval $[0,T]\subset[0,T_*)$ of the selected original
Schwartz-data branch,
\[
 V\in W^{1,\infty}([0,T];L^2),\qquad
 \|V_t(t)\|_2^2\le\frac94\int\rho(t)|F(t)|^2
 \quad\hbox{for almost every }t,
 \qquad F=\nu\Delta u-(u\cdot\nabla)u.
\]
Also $V\in L^\infty([0,T];H^1)$ and
\begin{equation}\label{st:drivers}
 \|V\|_6^2\le a_0D,\qquad
 \|V_t\|_2\le\tfrac32\|w\|_3^{1/2}\|F\|_3,
 \qquad a_0=9C_S^2/8.
\end{equation}
No bound uniform as $T\uparrow T_*$ is asserted.
\end{theorem}
\begin{proof}
The local package gives $u\in C^1([0,T];L^3)$ and a bounded $L^3$
norm. The preceding theorem makes $V$ Lipschitz into the separable
Hilbert space $L^2$, hence absolutely continuous with essentially
bounded strong derivative. One can obtain this last fact directly by
expanding in an orthonormal basis: the scalar coordinates are Lipschitz,
finite sums of their squared derivatives are bounded by the squared
Lipschitz constant, and the coordinate fundamental theorems assemble
into the Bochner integral and its almost-everywhere derivative.

In \eqref{st:difference} use the two times $t,t+h$ and the gradient
$g=\int_t^{t+h}\nabla p(s)\,ds\in\mathcal G_3$. The inclusion follows
from the earlier gradient decomposition and the local $C_tL^3_x$
regularity of $\nabla p$. The equation gives
$u(t+h)-u(t)+g=\int_t^{t+h}F(s)\,ds$. Divide by $h^2$ and let
$h\to0$ at a time of strong differentiability of $V$.
The representative is continuous in $L^3$, $F$ is continuous in $L^3$,
and their products converge in $L^1$ by H\"older. The right side tends
to $(9/4)\int\rho|F|^2$, proving the derivative bound.
This uses an admissible gradient increment, not a claim that the
pressure term is small.

The accepted weighted identity gives
$\|\nabla V\|_2^2\le9D/8$, while $\|V\|_2^2=3Q$.
Moreover $D=-\langle\rho w,\Delta u\rangle$ is continuous and bounded
on $[0,T]$, since $w$ is continuous in $L^3$ and $\Delta u$ in $L^3$.
Thus $V$ is locally bounded in $H^1$. Sobolev and H\"older give
\eqref{st:drivers}. No time derivative of $w$ at its zero set is used.
\end{proof}

\subsection{Finite smooth features instead of a moving sharp projection}

Let $\Pi$ be the orthogonal projection of matrices onto symmetric
trace-free matrices and let $\mathcal T M=\sum_{ij}R_iR_jM_{ij}$.
Write $C_p$ for a fixed bound of $\mathcal T:L^p(\mathbb R^3;
\mathbb R^{3\times3})\to L^p(\mathbb R^3)$, for $p=3/2,3$.
Use the same signed-work scalar as in Section~\ref{ss:section}:
\[
 \chi=\mathcal T\Pi(V\otimes V).
\]
Fix real $g_1,\ldots,g_n\in C_c^\infty((0,\infty))$ and $\eta>0$.
Set $\phi_j=g_j(\rho)$, $G_{ij}=\langle\phi_i,\phi_j\rangle$,
$b_i=\langle\chi,\phi_i\rangle$, and
\begin{equation}\label{st:ridge-definition}
 a=(G+\eta I)^{-1}b,\qquad
 R_{n,\eta}=\|\chi\|_2^2-b^{\mathsf T}(G+\eta I)^{-1}b,
 \qquad r=\chi-\sum_j a_j\phi_j.
\end{equation}
The letter $r$ here denotes a residual, not the speed $\rho$.

\begin{theorem}[Differentiable residual and exact signed pairing]
\label{st:residual}
All quantities in \eqref{st:ridge-definition} are locally absolutely
continuous in time; for $r$ this assertion is only needed and asserted
in $L^{3/2}$ on bounded spatial sets. The scalar residual satisfies
\begin{align}
 N_\rho^2\le R_{n,\eta}
 &=\|r\|_2^2+\eta|a|^2\le\|\chi\|_2^2,\label{st:bounds}\\
 K_u&=\langle\sigma,r\rangle,\qquad
 |K_u|\le\|\sigma\|_2\sqrt{R_{n,\eta}},\label{st:work}\\
 R_{n,\eta}'
 &=2\langle\chi,\chi_t\rangle-2a\cdot b'+a^{\mathsf T}G'a
  =2\langle r,\chi_t\rangle-2\left\langle r,\sum_j a_j\phi_{j,t}\right\rangle
 \quad\hbox{a.e.}\label{st:identity}
\end{align}
The first pairing in the last expression is an $L^3$--$L^{3/2}$
pairing; the second is an $L^2$ pairing. In particular they need not
be interpreted as an $L^2$ derivative of $\chi$.
\end{theorem}
\begin{proof}
On a compact classical interval write $M=\|V\|_6$ and $W=\|V_t\|_2$;
both are essentially bounded by Theorem~\ref{st:time}. Products of
time difference quotients give, distributionally in time,
\[
 \chi_t=\mathcal T\Pi(V_t\otimes V+V\otimes V_t),\qquad
 \|\chi\|_3\le C_3\sqrt{2/3}\,M^2,\quad
 \|\chi_t\|_{3/2}\le2C_{3/2}MW.
\]
Here the product rule follows by time mollification, or first on a
bounded spatial set by the pointwise absolutely continuous paths and
then by H\"older. These bounds also show that $\chi\in
W^{1,\infty}_tL^{3/2}_x\cap L^\infty_tL^3_x$; its membership in
$L^{3/2}$ follows separately from $V\in L^3$.
Furthermore $V\in C_tL^2_x\cap L^\infty_tL^6_x$ implies
$V\in C_tL^p_x$ for $2\le p<6$. In particular $\chi\in C_tL^2_x$.
Time mollification and H\"older in the dual spaces $L^3,L^{3/2}$ now
prove
$(\|\chi\|_2^2)'=2\langle\chi,\chi_t\rangle$ in $L^1$.
For clarity, mollification converges strongly in $L^q_tL^3_x$ for
finite $q$ on interior compact intervals, and the derivatives converge
strongly in $L^q_tL^{3/2}_x$; hence the product converges in $L^1$.
The continuous $L^2$ representative identifies the endpoint values.

The functions $h_j(z)=g_j(|z|^{2/3})$ are smooth and globally Lipschitz
on $\mathbb R^3$, being identically zero near zero and outside a ball.
Let $\ell_j=\sup|Dh_j|<\infty$. Since $\phi_j=h_j(V)$, the ordinary
chain rule on almost every pointwise time path, followed by Fubini, gives
\[
 \phi_{j,t}=Dh_j(V)V_t,\qquad
 \|\phi_{j,t}\|_2\le\ell_j W.
\]
These functions are in $C_tL^2_x\cap C_tL^3_x$. To check spatial
integrability, if the support of $g_j$ has lower endpoint $k_j>0$,
then $|\{\phi_j\ne0\}|\le k_j^{-3}\|w\|_3^3$; continuity in $L^3$
also follows from the uniform $L^\infty$ bound and continuity in $L^2$.
Consequently the scalar product rules are legitimate:
\[
 G_{ij}'=\langle\phi_{i,t},\phi_j\rangle+
          \langle\phi_i,\phi_{j,t}\rangle,\qquad
 b_i'=\langle\chi_t,\phi_i\rangle+\langle\chi,\phi_{i,t}\rangle.
\]
One can justify both by the same time mollification: the first term
in $b_i'$ uses $L^{3/2},L^3$ and the second uses $L^2,L^2$.
Since $G+\eta I\ge\eta I$, its inverse is absolutely continuous and
its derivative is $-(G+\eta I)^{-1}G'(G+\eta I)^{-1}$.
Differentiate the scalar expression for $R_{n,\eta}$ to obtain the
first identity in \eqref{st:identity}. Substitute the formulas for
$b'$ and $G'$ to get the second. All bounds are uniform on the chosen
compact interval for fixed $n,\eta$.

Completing the square in
$\|\chi-\sum d_j\phi_j\|_2^2+\eta|d|^2$ proves that its minimum is
$R_{n,\eta}$ at $d=a$. It also proves the middle identity and upper
bound in \eqref{st:bounds}. Its lower bound follows because each
$\phi_j$ belongs to $\mathcal H_\rho$. The shell cancellation of
Theorem~\ref{ss:shell} annihilates each $\phi_j$, proving
\eqref{st:work} from the previously established signed pairing.
The stated local absolute continuity of $r$ follows from that of
$\chi$, $a$, and the features, using the inclusion of $L^2$ into
$L^{3/2}$ on bounded spatial sets.
\end{proof}

\begin{corollary}[An explicit residual evolution bound]\label{st:evolution}
Let $A_3=(\sum_j\|\phi_j\|_3^2)^{1/2}$ and
$\ell=(\sum_j\ell_j^2)^{1/2}$. For $R=R_{n,\eta}$ one has a.e.
\begin{equation}\label{st:evolution-bound}
 |R'|\le 4C_{3/2}MW
       \left(C_3\sqrt{2/3}\,M^2+A_3\sqrt{R/\eta}\right)
       +2\ell\eta^{-1/2}WR.
\end{equation}
For each fixed snapshot there is a single nested choice of smooth
feature spans, independent of the snapshot, for which
$R_{n,\eta_n}\downarrow N_\rho^2$ whenever $\eta_n\downarrow0$.
\end{corollary}
\begin{proof}
The identity \eqref{st:bounds} gives $\|r\|_2\le\sqrt R$ and
$|a|\le\sqrt{R/\eta}$. Hence
$\|r\|_3\le C_3\sqrt{2/3}M^2+A_3\sqrt{R/\eta}$ and
$\|\sum a_j\phi_{j,t}\|_2\le\ell W\sqrt{R/\eta}$.
Apply the two indicated H\"older inequalities to \eqref{st:identity}
and use the estimate for $\chi_t$.

Choose a countable family in $C_c^\infty((0,\infty))$ whose real
linear span is uniformly dense in continuous compactly supported
functions on each compact speed interval, using rational breakpoints,
rational coefficients and smooth rational-width mollifications.
For any fixed speed field, composition is dense in $\mathcal H_\rho$:
on a compact speed interval the pushforward measure is finite, continuous
functions are dense in its $L^2$, and exhaustion handles both speed tails.
Enumerate this family and take its first $n$ elements. Enlarging the
span and decreasing the positive penalty makes the minimum nonincreasing.
For any fixed finite linear combination approximating $P_\rho\chi$,
its coefficient penalty tends to zero as $n\to\infty$.
Comparison with this competitor and then approximation prove that the
limit is at most $N_\rho^2$; \eqref{st:bounds} gives the reverse bound.
This is a pointwise-in-time convergence statement, not convergence of
the derivatives or an estimate uniform at a singular endpoint.
\end{proof}

\paragraph{What this unblocks, and the remaining bound.}
For nonzero $w$ set
$\delta_{n,\eta}=3\sqrt{R_{n,\eta}}/(2\|\rho^3\|_2)$, with value
zero for $w=0$. Then $\delta_\rho\le\delta_{n,\eta}\le1$, and
\eqref{st:work} supplies the same fourth-power inequality as
\eqref{ss:depleted-diff} with $\delta_{n,\eta}$ in place of
$\delta_\rho$. The new content is the justified evolution
\eqref{st:identity}--\eqref{st:evolution-bound}, not a new assumed
continuation theorem. It does not differentiate $P_\rho$ or select an
optimizer over all speed functions.

The first unsupported step in an attempted closure is an input-only
endpoint bound for the driver in \eqref{st:evolution-bound}:
$M^2\le a_0D$ and
$W^2\le(9/4)\int\rho|\nu\Delta u-(u\cdot\nabla)u|^2$ are local
identities or bounds, not energy-controlled endpoint integrals. In
addition the displayed derivative estimate deteriorates as $\eta\downarrow0$
and as more features with large Lipschitz constants are included.
Pointwise monotone convergence of the residuals does not remove these
losses. No bound on $L_c$, on $\int\|\sigma\|_2^4$, or on these drivers
in terms of arbitrary initial data is obtained. Thus the temporal
regularity calculation is complete at fixed regularization, but the
arbitrary-data critical producer and NS-R3 remain open.
